\section{Supporting analysis -- Giải thích các thành phần}
\subsection{Cache trên cùng output và provider tĩnh}
Current giữ packet hiện tại; epoch60 giữ metadata trong đoạn 60 s; session tích lũy cả phiên.
Các chế độ dùng cùng GPS/reference, không tạo Q mới. Đây là phép tách ảnh hưởng của bộ nhớ khỏi
ảnh hưởng của privacy kernel.
\Needspace{210pt}
\par\begingroup\small\setstretch{1.04}\setlength{\parskip}{0pt}
\begin{longtable}{@{}P{4cm}P{3.719cm}P{3.719cm}P{3.719cm}@{}}
\caption{Macro Recall@5 theo phạm vi cache trên cùng output}\label{tab:cache}\\
\toprule
\textbf{Phương pháp} & \textbf{Current (\%)} & \textbf{Epoch60 (\%)} & \textbf{Session (\%)} \\
\midrule\endfirsthead
\multicolumn{4}{l}{\small\textit{Bảng \thetable\ (tiếp)}}\\\toprule
\textbf{Phương pháp} & \textbf{Current (\%)} & \textbf{Epoch60 (\%)} & \textbf{Session (\%)} \\
\midrule\endhead
\midrule\multicolumn{4}{r}{\small\textit{Còn tiếp ở trang sau}}\\\endfoot
\bottomrule\endlastfoot
Raw GPS & 99{,}999 & 99{,}999 & 99{,}999 \\
Geo-I/REM Z (K1) & 98{,}610 & 98{,}610 & 99{,}700 \\
TransProtect* & 99{,}996 & 99{,}996 & 99{,}996 \\
Semantic correlation* & 99{,}999 & 99{,}999 & 99{,}999 \\
Fake-query insertion* & 99{,}999 & 99{,}999 & 99{,}999 \\
\textbf{GeoI-BQ v1} & 96{,}780 & 97{,}123 & 99{,}595 \\
\end{longtable}\endgroup\par


GeoI-BQ tăng từ 96,780\% current lên 97,123\% epoch60 và 99,595\% session. Cache giúp bù việc
không truy vấn trực tiếp GPS. Tuy nhiên, proof inclusion chỉ giữ ở cùng GPS/request và metadata
hợp lệ. Di chuyển sang một vùng khác có thể làm reference đổi; thông tin giờ mở cửa/giá/status
động có thể làm cache mất hiệu lực. Do đó kết quả session tĩnh không là chứng chỉ chất lượng
dịch vụ online động mà không có thêm chính sách expiry/refresh.

\subsection{Response depth và vai trò của GPS local}
Trên cùng Q/time và provider signatures, truncation L30 cho macro Recall 94,170\%, thấp hơn L60
5,425 điểm phần trăm, CI95\% [3,486;7,527]. Radius chịu chênh lệch lớn nhất, khoảng 10,374 điểm
phần trăm. Bảng chi tiết ở Phụ lục~\ref{app:support}; đây là replay/sensitivity có contract riêng,
không là một lần cải thiện privacy hoặc một thử nghiệm thành phố mới.

Một correction metadata đã đổi nhãn request nguồn cũ từ L30 sang L60 để khớp callback thực tế.
Nó không thay Q/time, reply masks, Recall hoặc độ dài bytes trong main trial. Cấu hình trình bày
hiện tại ghi L60 mỗi category/tọa độ; không nối score của L30 cũ với score mới thành một comparison
cùng cấu hình. Source hashes và verifier ghi lại correction này.

GPS local dùng đánh giá/ranking tại grid 20 s; lớp bảo vệ chỉ đưa private information mới vào
upstream theo pacing tối thiểu 60 s và cap. Vì vậy ít protected reads không đồng nghĩa điện thoại
chỉ bật GNSS từng ấy lần, cũng không chứng minh tiết kiệm pin. Đây là ranh giới giữa quan sát
của server và nhu cầu local, không phải một thí nghiệm năng lượng.

\subsection{Ablation public base và thích nghi theo belief}
Bảng ablation nội bộ ngày 14/10 giữ cùng frozen release ở Phụ lục~\ref{app:support}. Matched public
base $\Gamma=0$ đạt macro 99,711\%, cao hơn GeoI-BQ $\Gamma=1$ là 99,595\% trên workload này.
Service-cover đạt 100\% nhưng là control retrieval nội bộ, không phải một paper đối chứng.

Kết quả đó phù hợp với định lý utility: gain score không bắt buộc actual Recall gain. $\Gamma=0$
loại private adaptation của Q khi context cố định; vì thế nó cũng có Q-channel independence mạnh
hơn. Cohort này \textbf{chưa chứng minh belief adaptation cần thiết hoặc $\Gamma=1$ tối ưu}.
Thiết kế đề xuất cho phép bounded adaptation có accounting, còn giá trị thực tế của adaptation
trong catalogue lớn/động cần được xác nhận. Không thay baseline sau khi nhìn TEST để che kết quả này.

\section{Discussion -- Độ tin cậy và phạm vi kết luận}\label{sec:discussion}
\subsection{Mối quan hệ giữa chứng minh và benchmark}
Lý thuyết trả lời một câu hỏi khác với accuracy của một attack: các phân phối transcript được
phép khác nhau tới đâu khi private input đổi. Nếu một attacker mới mạnh hơn làm Hit100 tăng,
định lý kernel không tự bị phủ định; phải kiểm tra liệu attack/side channel có nằm trong model và
có vi phạm cận phân phối đã chứng minh không. Ngược lại, Hit thấp không đủ chứng nhận implementation
đúng kernel lý tưởng. Hai loại bằng chứng phải được đọc cùng giả thiết.

Với S1--S3/S9--S10, finite bank hiện tại cho thấy Q-only hạn chế định vị tốt hơn các controls/adaptations
trong cùng trial. Với S4/S6, có thêm bằng chứng tác vụ có target cụ thể. S5 chưa có positive control
đủ mạnh; S7 là conditional interface property chung; S8 là diagnostic tương quan tổng hợp và raw side
information. Không lấy việc đã có mười dòng scenario làm bằng chứng mười bài toán đều được giải quyết đầy đủ.

\subsection{Validity của phép so sánh phương pháp}
Các adaptation giữ ý tưởng cơ chế nhưng thay neural predictors hoặc nguồn semantic. Do đó conclusion
chỉ là controlled comparison của các implementation đã audit. Epsilon, output contracts và query
counts khác nhau; không nói cùng privacy budget hoặc cùng cost. Native metrics N/A phản ánh thiếu
contract tương ứng, không loại một result bất lợi. Cần author code, dữ liệu và task alignment trước
khi phát biểu superiority đối với kết quả chính thức của paper.

Raw control xác nhận năm coordinate targets có thể quan sát, nhưng S5 cho thấy control có thể không
cung cấp tín hiệu trong một cutoff cụ thể. S8 V2 cho thấy một bank học có thể thiếu fusion rất đơn giản.
Vì vậy attacker selection phải công khai, giữ loss-specific models và không trình bày số lớn ở MAE
như một chứng chỉ attacker tối ưu đã thất bại.

\subsection{Validity của dữ liệu và utility gần bão hòa}
Chia family tránh một dạng leakage giữa các chuyến cùng nhóm, nhưng bản đồ/cohort đã được xem trong
quá trình phát triển. TRAIN/SELECTION/TEST không biến nó thành một geography holdout độc lập. Native
SUMO phản ánh những ràng buộc giao thông của mạng đã xây, chưa đại diện hành vi người thật, quyết định
POI hay nhãn người/xe. Bootstrap family/block đo resampling của số nhóm hữu hạn đang có, không đo
tất cả nguồn domain shift.

Catalogue có 418 POI và sáu categories. K5 nhân L60 mỗi category có thể trả tối đa 1.800 entries trước
dedup ở một packet, nên coverage và session cache dễ gần bão hòa. Đây là lời giải thích quan trọng
cho Recall gần 100\%, không chỉ là đặc tính của planner. Nếu API cho tải toàn catalogue hợp lệ một
lần với cost chấp nhận được, client có thể giảm nhu cầu periodic coordinate queries; lợi thế của
planner khi đó phải được đánh giá lại theo chính contract đó. Không ghép số từ một diagnostic bulk
khác cohort/cấu hình vào bảng chính để chứng minh ưu thế.

Privacy MAE lớn cũng phụ thuộc public library rộng so với route workload. Công bố sự khác nhau này
giúp phân biệt guard lý thuyết với lợi ích do geometry hiện tại. Một nghiên cứu tiếp theo cần catalogue
lớn/động hơn, reference đầy đủ và những family mới chốt trước khi đánh giá.

\subsection{Giới hạn triển khai và tài nguyên}
Float Geo-I, entropy nguồn random và full-state restore là những điều kiện thực thi còn phải kiểm
chứng trước deployment có cam kết. Header/account/IP, kích thước payload không được giữ cố định,
click feedback và private start/stop timing có thể tạo kênh khác. Finite cap của một phiên không
xóa composition qua nhiều phiên liên kết hoặc qua người tương quan.

Chi phí hiện báo là counts và serialized index JSON. Metadata thực, network latency, server billing,
GNSS duty cycle và năng lượng chưa có phép đo đồng bộ. BQ/REM-Z được replay nên không có generation
timing cùng trial. Các giới hạn này không làm mất giá trị của counts đã kiểm chứng, nhưng ngăn kết
luận phương pháp nhanh hơn/rẻ hơn trên thiết bị người dùng.

\subsection{Kiểm chứng tái lập và failure accounting}
Main verifier kiểm tra độc lập 392.832 ô utility, 600 family-level privacy scores, 528 native-cost
trip values, 5.456 proposed prefixes và tổng 104.168 point queries trên toàn các split/arms.
Supplement kiểm tra split/block, future cutoffs, selectors, predictions, CIs và mọi fake packet.
Kết quả failure gate S8 ban đầu được giữ bên cạnh checker sửa điều kiện kiểm tra, không ghi đè
predictions. Cả ba benchmark folders và 108 frozen release pins được xác thực trước xuất bảng.

Quy trình này không biến dữ liệu đã xem thành dữ liệu mới hoặc làm mọi sai số biến mất. Nó bảo đảm
người review có thể đi từ kết luận về đúng con số, task, output và source dùng để tính. Trong luận
văn, đây là cơ sở để giới hạn contribution ở một phương pháp có đặc tả rõ, chuỗi chứng minh có điều
kiện và một đánh giá thực nghiệm có thể kiểm tra.
